\documentclass[11pt]{article}
\usepackage[margin=1in]{geometry}
\usepackage{booktabs}
\usepackage{amsmath}
\usepackage{microtype}
\usepackage{xcolor}
\usepackage[colorlinks=true,linkcolor=black,urlcolor=blue]{hyperref}
\usepackage{caption}
\captionsetup{font=small,labelfont=bf}

\newcommand{\pos}[1]{\textcolor{blue!60!black}{#1}}
\newcommand{\bad}[1]{\textcolor{red!65!black}{#1}}

\title{\vspace{-2.5em}Two Ways to Lose a Fantasy Football League}
\author{An empirical decomposition of Fat Man's Fantasy, 2016--2025}
\date{September 2026}

\begin{document}
\maketitle
\vspace{-1.5em}

\begin{abstract}
\noindent Across 56 manager-seasons we decompose league performance into the
three components a manager actually controls: the draft, in-season roster
acquisition, and weekly lineup selection. Two managers finish near the bottom of
the league by mutually exclusive routes. \texttt{Gordonulus} is an above-average
in-season manager who destroys his rosters in the first four rounds.
\texttt{mikeion} drafts adequately and then stops playing. We also find no
evidence that drafting ability persists across eras $(r = -0.157)$, which limits
how much of this should be read as skill rather than circumstance.
\end{abstract}

\section{Framework}

A season's starting-lineup points can be written as
\[
  \text{Points} \;=\; \underbrace{f(\text{draft})}_{\text{August}}
  \;+\; \underbrace{g(\text{acquisitions})}_{\text{September--December}}
  \;-\; \underbrace{h(\text{lineup error})}_{\text{every Sunday}} \;+\; \varepsilon .
\]
Each term is measured on a common scale. A player's contribution is his half-PPR
points in weeks 1--14 (playoff weeks are excluded: they decide nothing for a team
already eliminated), converted to points above replacement at his own position,
with replacement recomputed each season for a 14-team lineup carrying two flex
spots. A draft pick is then scored against what a pick at that slot typically
returns, so that picking first is not mistaken for drafting well.

Regressing season wins on the three components gives

\begin{table}[h]
\centering
\small
\begin{tabular}{lrrr}
\toprule
Component & Wins per SD & Raw $r$ & Partial $r$ \\
\midrule
Draft value            & $+1.19$ & $+0.52$ & --- \\
Lineup efficiency      & $+0.54$ & $+0.33$ & $+0.34$ \\
Waiver production      & $+0.33$ & $+0.01$ & $+0.25$ \\
\bottomrule
\end{tabular}
\caption{$n = 56$ manager-seasons, $R^2 = 0.378$.}
\end{table}

\noindent Waiver production correlates with wins at essentially zero until the
draft is held constant, at which point it recovers to $+0.25$. The two effects
had been cancelling: a failed draft \emph{forces} a manager onto the wire, so raw
waiver points measure draft failure as much as waiver skill. Note also that
$R^2 = 0.378$. Roughly $62\%$ of a season is schedule, injury, and luck.

\section{The two failure modes}

Six controllable quantities, $z$-scored across the fourteen managers with at
least two full seasons. Higher is better throughout.

\begin{table}[h]
\centering
\small
\begin{tabular}{lrrrr}
\toprule
& \multicolumn{2}{c}{$z$-score} & \multicolumn{2}{c}{raw value} \\
\cmidrule(lr){2-3}\cmidrule(lr){4-5}
Dimension & \texttt{mikeion} & \texttt{Gordonulus} & \texttt{mikeion} & \texttt{Gordonulus} \\
\midrule
Waiver volume (claims/season) & \bad{$-1.74$} & \pos{$+0.19$} & $10.3$ & $32.5$ \\
Bid discipline (\$ over 2nd bid) & \bad{$-0.90$} & \bad{$-0.86$} & $\$7.44$ & $\$7.33$ \\
Waiver efficiency (points/\$) & \bad{$-1.83$} & \bad{$-0.65$} & $4.1$ & $6.4$ \\
Lineup setting (\% of optimal) & \bad{$-0.74$} & \bad{$-0.96$} & $91.8\%$ & $91.6\%$ \\
Early-round value (per pick) & \bad{$-1.23$} & \bad{$-1.48$} & $-26.1$ & $-31.4$ \\
Bust avoidance (\% of R1--4) & \bad{$-1.45$} & \bad{$-2.03$} & $31\%$ & $38\%$ \\
\bottomrule
\end{tabular}
\caption{League means: $30.3$ claims, $\$5.21$ overpay, $7.6$ points/\$,
$92.5\%$ optimal, $-0.6$ per early pick, $16\%$ bust rate.}
\end{table}

\subsection*{\texttt{Gordonulus}: a draft problem}

His waiver volume is the only positive number either manager records. He is an
active, engaged in-season manager --- $32.5$ claims a season against a league
mean of $30.3$, and he twice acquired Matthew Stafford for \$3, returning $302$
points for six dollars across two seasons.

He loses in August. His bust rate in rounds 1--4 is $38\%$ against a league rate
of $16\%$, the worst mark on any dimension recorded by either manager. His
early-round picks average $-31.4$ points of value each; over four picks that is
roughly $125$ points a season surrendered before Week 1. His rounds 8--14 are
also below water at $-11.5$ per pick, so there is no late-round recovery.

The pattern is not obviously about age or position --- his round 1--4 skill picks
average $25.0$ years old, \emph{younger} than the league's $26.3$. Nick Chubb at
pick 6 in 2023 $(-157.3)$ and Elijah Mitchell at pick 47 in 2022 $(-93.8)$ were
both bets on running backs returning from injury.

\subsection*{\texttt{mikeion}: a participation problem}

His draft is mediocre rather than catastrophic, and it is specifically
front-loaded: $-26.1$ per pick in rounds 1--4 but $+9.1$ per pick from round 8
onward, one of the better late-round records in the league. He is not bad at
evaluating players; he is bad at the four picks that matter most.

What separates him from everyone else is the season that follows. At $10.3$
claims per season he is last in the league by a factor of three, and he converts
$4.1$ points per FAAB dollar against a league mean of $7.6$. In absolute terms he
extracts $252$ points a season from the waiver wire; the league median is $692$
and the leader is $1{,}098$. Over a fourteen-week season that gap is roughly $60$
points a week of unclaimed production.

His claim \emph{win rate} of $63\%$ is the second highest in the league, which is
diagnostic rather than flattering. Across the fourteen managers,

\[
  r(\text{claims per season},\ \text{waiver points}) = +0.888, \qquad
  r(\text{price paid per win},\ \text{waiver points}) = -0.799,
\]
\[
  r(\text{claim win rate},\ \text{waiver points}) = -0.369 .
\]

\noindent A high win rate means bidding only when certain, and certainty is
expensive: he pays \$9.40 per successful claim against a league mean of \$5.50.
Volume, not accuracy, is what the wire rewards.

\section{What this is not}

Three caveats materially limit the above.

\paragraph{Drafting ability does not persist.} Extending the analysis to the
MyFantasyLeague seasons of 2016--17, the correlation between a manager's draft
performance then and now is $r = -0.157$, $95\%$ CI $[-0.67, +0.46]$. Eight of
twelve managers changed sign. The manager with the best four-year draft record in
the current era had the worst in the previous one. Pooling all six gradeable
seasons, the share of variance attributable to the manager falls from $0.215$ to
$0.027$.

\paragraph{The early-round samples are small.} Each manager has 16 picks in
rounds 1--4 across four seasons. A $38\%$ bust rate is six events. The waiver
figures rest on 1{,}089 auctions and are considerably more solid.

\paragraph{The dimensions are not equally weighted.} A $z$-score conceals effect
size. The entire league sits between $91.1\%$ and $94.1\%$ lineup efficiency; one
standard deviation is about $2.8$ points per week. The waiver and early-round
columns carry real magnitude, the lineup column does not.

\section{Conclusion}

Both managers finish near the bottom, and the remedies do not overlap. One should
draft more conservatively in the first four rounds and would otherwise be fine.
The other should open the application on Tuesdays. Neither is doing the thing the
other needs to fix, which is why the shared complaint --- that they are simply
unlucky --- is only $62\%$ correct.

\end{document}
